% Source-derived chapter 05: Strange duality
% Original source: rigid-local-systems-multiplicative-eigenvalue-problem
\chapter{Strange duality}
\label{rigid-local-systems-multiplicative-eigenvalue-problem-chapter-05}
\source{rigid-local-systems-multiplicative-eigenvalue-problem}

\begin{remark}[Source section]
\label{rigid-local-systems-multiplicative-eigenvalue-problem-chapter-05-source}
\source{rigid-local-systems-multiplicative-eigenvalue-problem}
\source{rigid-local-systems-multiplicative-eigenvalue-problem}
This chapter reproduces the corresponding section of the retrieved
LaTeX source, including its definitions, statements, examples, and
proofs. It has not yet been translated into Lean.
\notready
\end{remark}

% The source section title was: Strange duality
\label{SD}
\source{rigid-local-systems-multiplicative-eigenvalue-problem}
\subsection{Proof of Theorem \ref{egregium}}
 We show how to  associate a level $\ell$ representation $\lambda$ of $\op{SL}(n)$ to  a conjugacy class in $\op{GL}(\ell,\Bbb{C})$ of order dividing $n$:
\begin{defi}
\source{rigid-local-systems-multiplicative-eigenvalue-problem}
\notready \label{reprem2}
\source{rigid-local-systems-multiplicative-eigenvalue-problem}
Suppose $\bar{A}\in \op{GL}(\ell,\Bbb{C})$ such that  $A^n=I_{\ell}$.   We can find an $\ell\times n$ Young diagram $\mu$, with $n> \mu_1\geq \mu_2\geq \dots \geq \mu_{\ell}\geq 0$, such that $\bar{A}$ is conjugate to the diagonal matrix with entries $\exp(2\pi\sqrt{-1}\mu_i/n)$, for $i=1,\dots,\ell$.

  Let $\lambda=\lambda(\bar{A})=\mu^T$, the transpose Young diagram of $\mu$, which fits in a box of size $n\times \ell$. Note that $\lambda$ is dominant integral weight of $\op{SL}(n)$ of level $\ell$, i.e., $\lambda_1-\lambda_n \leq \ell$, and also  $\lambda_n=0$, which is  normalisation condition on $SL(n)$ representations (see Sections \ref{loudvoice} and \ref{loudvoice2}).
\end{defi}


\begin{lemma}
\source{rigid-local-systems-multiplicative-eigenvalue-problem}
\notready\label{reprem}Let $\ell$ and $n$ be fixed. The correspondence $\bar{A}\to \lambda(\bar{A})$ in Definition \ref{reprem2} gives a bijective correspondence between the  set of semisimple conjugacy classes  $\bar{A}\in\op{GL}(\ell,\Bbb{C})$ with $\bar{A}^n=I_{\ell}$,  and the set of  level $\ell$ representations of $\op{SL}(n)$.
\source{rigid-local-systems-multiplicative-eigenvalue-problem}
\end{lemma}
\begin{proof}
The reverse correspondence is the following: A  level $\ell$ representation of $\op{SL}(n)$
gives a Young diagram  $\lambda$ that fits in an $n\times \ell$ box with $\lambda_n=0$ (the  normalisation condition).   Let $\mu=\lambda^T$.  Clearly $\mu_1<n$ since $\lambda_n=0$. The conjugacy class corresponding to $\lambda$ is then the diagonal matrix with entries $\exp(2\pi\sqrt{-1}\mu_i/n)$, for $i=1,\dots,\ell$.
\end{proof}


Let $\ell$ and $n$ be fixed. Applying Lemma \ref{reprem}, coordinate by coordinate, we find that the following sets are in bijective correspondence:
\begin{enumerate}
\item[(A)]The set of $s$-tuples of semisimple  conjugacy classes $\mathcal{A}=(\bar{A}_1,\dots,\bar{A}_s)$ in $\op{GL}(\ell,\Bbb{C})$, with $\bar{A}_i^n=1$, such that  $\prod_{i=1}^s \det \bar{A}_i =1$.
\item[(B)] The set of grade zero bundles $\mathcal{B}(\vec{\lambda},\ell)$  of level $\ell$ on $\Par_{n,\mathcal{O},S}$, i.e., $\vec{\lambda}$ is a tuple of  dominant integral weights  for $\op{SL}(n)$ of level $\ell$, and  $n$ divides $\sum_{i=1}^s|\lambda^i|$.
\end{enumerate}
The correspondence between (A) and (B) is 
\begin{equation}\label{corr2}
\source{rigid-local-systems-multiplicative-eigenvalue-problem}
\mathcal{A}\mapsto (\lambda^1,\dots,\lambda^s),\ \ \lambda^i=\lambda(\bar{A}_i), \ i=1,\dots,s
\end{equation}
(also see Remark \ref{newRem} for the compatibility with the point $v(\mathcal{A})$ from Definition \ref{finiteness}). 

The condition on  the grade in  (B) is that $n$ divides $\sum_{i=1}^s|\lambda^i|$. Since $|\mu|=|\mu^T|$ for a Young diagram $\mu$, this is the same condition, under the correspondence, as the determinant condition in (A).
\begin{remark}
\source{rigid-local-systems-multiplicative-eigenvalue-problem}
\notready\label{newRem}
\source{rigid-local-systems-multiplicative-eigenvalue-problem}
In the above correspondence \eqref{corr2}, the point $v(\mathcal{A})\in \Delta_n^s$ from Definition \ref{finiteness} is
$(\kappa(\lambda^1/\ell),\dots,\kappa(\lambda^s/\ell))$ (see the Killing form isomorphism $\kappa$ from Section \ref{loudvoice}, and Theorem \ref{blacktea}).

\end{remark}

\iffalse
As in Section \ref{newRigid} assume  that the  eigenvalues of $\bar{A}_i$ are $\exp(2\pi\sqrt{-1}\mu_j^i/n)$ with integers $0\leq \mu_j^i<n$, so that the $\mu^i=(\mu^i_1,\dots,\mu^i_{\ell})$ are Young diagrams that fit in an $\ell\times n$ box, i.e.,  for $i\in [s], n> \mu^i_1\geq \mu^i_2\geq \dots \geq \mu^i_{\ell}\geq 0$. This gives a parametrization of (A). we also have $\sum_{i,a}\mu^i_a$ divisible by $n$

The line bundles in (B) are of the form $\mathcal{B}(\vec{\lambda},\ell)$, here $\lambda^j$ are dominant integral weights  of level $\ell$, $j=1,\dots,s$, i.e.,
\begin{equation}\label{weyl}
\source{rigid-local-systems-multiplicative-eigenvalue-problem}
\lambda^{j}_1\geq \lambda^{j}_2\geq \dots\geq \lambda^{j}_n\geq \lambda^{j}_1-\ell.
\end{equation}
Note that we may assume that $\lambda^j_n=0$ for all $j$ because we are considering weights for $SL(n)$ here.

\end{proof}
 Then
the correspondence takes $\mathcal{A}$ to $\mathcal{B}(\vec{\lambda},\ell)$ with $\lambda^i=(\mu^i)^T$.
Note that $\lambda^i_n=0$, so the weights of $\lambda$ are normalised for $\op{SL}_n$. The grade zero condition in (B) is equivalent to the product of determinants condition in (A).
\fi
The following numerical strange duality \eqref{GepnerPP} is a direct consequence of results of Witten and Gepner \cite{Gepner,Witten}. The full strange duality asserts that the vector spaces in \eqref{GepnerPP} are in fact  canonically dual to each other (see Section \ref{understanding}), but we will not need this fact.
\begin{proposition}
\source{rigid-local-systems-multiplicative-eigenvalue-problem}
\notready\label{BWag}
\source{rigid-local-systems-multiplicative-eigenvalue-problem}
Suppose $\ml=\mathcal{B}(\vec{\lambda},\ell)$ is grade zero line bundle on $\Par_{n,\mathcal{O},S}$ with $\ell>0$, $\vec{\lambda}$ an $s$-tuple of dominant integral weights, which are normalised, i.e., $\lambda^i_n=0, i\in[s]$. Write $\sum_{i=1}^s|\lambda^i|=nu$. Let  $\widetilde{\mn}$ be a line bundle of degree $-u$ on $\Bbb{P}^1$. Let $\ml'$ be the grade zero line bundle $\mathcal{B}(\vec{\lambda}^T,n)$ on $\Par_{\ell,\widetilde{\mathcal{N}},S}$. Then,
\begin{equation}\label{GepnerPP}
\source{rigid-local-systems-multiplicative-eigenvalue-problem}
h^0(\Par_{n,\mathcal{O},S},\ml)=h^0(\Par_{\ell,\widetilde{\mathcal{N}},S},\ml')
\end{equation}
\end{proposition}
\begin{proof}
The rank of $H^0(\Par_{n,\mathcal{O},S},\ml)$ equals a generalized Gromov-Witten number $\langle\sigma_{I^1}, \sigma_{I^2},\dots,\sigma_{I^s}\rangle_{0,D}$ with $D=\ell-u$, $I^1,\dots,I^s$ are subsets of $[n+\ell]$ each of cardinality $n$ corresponding to the partitions $\vec{\lambda}$ (see Definition  \ref{correspondence} and Lemma \ref{newlabel2}).

Now interpret the GW count as a count of subbundles of rank $\ell$ of degree $D$ of a generic bundle of degree $D$ by forming the (ordinary) dual of the quotient $(\mw/\mv)^*\subseteq \mw^*$ above. Using \eqref{newlabel3}, we  write the numbers above as $h^0(\Par_{\ell,\mathcal{N}',S},\mathcal{B}(\vec{\mu},n))$, with $\deg(\mathcal{N}')=D=\ell-u$. Using Lemma \ref{threepointone}, these numbers are equal to $h^0(\Par_{\ell,\widetilde{\mathcal{N}},S},\ml')$, which finishes the proof.


%\begin{equation}\label{GepnerP}
%h^0(\Par_{n,\mathcal{N},S},\ml)=\langle\sigma_{I^1}, \sigma_{I^2},\dots,\sigma_{I^s}\rangle_{d,D}=h^0(\Par_{\ell,\widetilde{\mathcal{N}},S},\ml')
%\end{equation}


\end{proof}

\begin{proposition}
\source{rigid-local-systems-multiplicative-eigenvalue-problem}
\notready\label{schubertt}
\source{rigid-local-systems-multiplicative-eigenvalue-problem}
The following are equivalent in the correspondence  between (A) and (B):
\begin{enumerate}
\item There is a unitary representation of $\pi_1$ with  local monodromy data $\mathcal{A}$.
\item The line bundle $\mathcal{B}(\vec{\lambda},\ell)$ on $\Par_{n,{\mathcal{O}},S}$ is effective.
\item $v(\ma)$ is a point of $P_n(s)$ (see Definition \ref{finiteness}).
\end{enumerate}

\end{proposition}
\begin{proof}
As in Prop \ref{BWag}, write $\sum|\mu^i|=nu$. Let $\widetilde{\mathcal{N}}$ be a line bundle of degree $-u$ on $\Bbb{P}^1$.
There is a unitary representation with  local monodromy data $\mathcal{A}$ if and only if $\ml'=\mathcal{B}(\vec{\mu},n)$ is effective on the moduli stack  $\Par_{\ell,\widetilde{\mathcal{N}},S}$ (use Corollary \ref{sansD} with $n$ and $\ell$ interchanged): We have an equation
$\sum|\mu^i| +n\deg \widetilde{\mathcal{N}} =0\cdot\ell$.
%Note when we apply strange duality $\deg{\mn}=n$. But $\Par_{n,{\mathcal{N}},S}$ can be identified with $\Par_{n,\mathcal{N}(-n),S}=\Par_{n,{\mathcal{O}},S}$.

By numerical strange duality, i.e., equality \eqref{GepnerPP} in Proposition \ref{BWag},  $\ml'$ is effective if and only if $\ml$ on  $\Par_{n,{\mathcal{O}},S}$ is effective. Since the grade of $\ml$ is zero, it is effective if and only if $(\kappa(\lambda^1/\ell),\dots,\kappa(\lambda^s/\ell))$ is a point of $P_n(s)$ (use Proposition \ref{kappa}).
It is easy to check that $v(\ma)=(\kappa(\lambda^1/\ell),\dots,\kappa(\lambda^s/\ell))$, which finishes the proof.
\end{proof}

\begin{defi}
\source{rigid-local-systems-multiplicative-eigenvalue-problem}
\notready\label{Shosta}
\source{rigid-local-systems-multiplicative-eigenvalue-problem}
In the setting of Prop \ref{schubertt}, we will call the moduli of unitary representation of $\pi_1$ with local monodromy data $\mathcal{A}$ (possibly empty)  ``strange dual'' to  the line bundle  $\mathcal{B}(\vec{\lambda},\ell)$ (possibly non-effective), i.e., the numerical data (A) and (B) are said to be in strange duality (also see Remark \ref{added}).

\end{defi}
\begin{defi}
\source{rigid-local-systems-multiplicative-eigenvalue-problem}
\notready\label{amirK}
\source{rigid-local-systems-multiplicative-eigenvalue-problem}
We will also assign parabolic weights if the $\lambda_i$ are dominant and $\ell>0$: The parabolic weights of $\ml=\mathcal{B}(\vec{\lambda},{\ell})$ are $\frac{{\lambda^i}}{\ell}$ at $p_i$. The parabolic degree of a point of $\Par_{n,\mathcal{O},S}$ corresponding to the line bundle  $\mathcal{B}(\vec{\lambda},{\ell})$ is $\frac{1}{n}(\deg{\mathcal{O}}+\frac{1}{\ell}\sum_{i=1}^s|\lambda^i|)=\ell-1-\deg\widetilde{\mn}$. The parabolic degree of $\mathcal{B}(\vec{\lambda}^T,n)$ on $\Par_{\ell,\widetilde{\mathcal{N}},S}$ in Proposition \ref{BWag} is then $\frac{1}{\ell}(\deg{\widetilde{\mn}}+\frac{1}{n}\sum_{i=1}^s|(\lambda^i)^T|) =n-n=0$.
\end{defi}

\begin{theorem}
\source{rigid-local-systems-multiplicative-eigenvalue-problem}
\notready\label{christmas} 
\source{rigid-local-systems-multiplicative-eigenvalue-problem}
Under the correspondence between the sets (A) and (B) above, suppose that $\mathcal{B}(\vec{\lambda},\ell)$ on $\Par_{n,{\mathcal{O}},S}$ is effective (equivalently, by Proposition \ref{schubertt}, there is a unitary representation of $\pi_1$ with  local monodromy data $\mathcal{A}$). Then
\begin{enumerate}
\item[(i)] $\mathcal{B}(\vec{\lambda},\ell)$ is a Hilbert basis element of  the semigroup $\Pic^+(\Par_{n,\mathcal{O},S})$ if and only if 
any unitary local system on $\Bbb{P}^1-S$ with local monodromy data $\mathcal{A}$ is  irreducible.
\item[(ii)] $\mathcal{B}(\vec{\lambda},\ell)$ is an F-line bundle on   $\Par_{n,\mathcal{O},S}$ (and hence gives rise to a F-vertex of $P_n(s)$ by Lemma \ref{schubertt}) if and only  there is a unitary, irreducible rigid local system on  $\Bbb{P}^1-S$ with local monodromy data $\mathcal{A}$. 
\end{enumerate}
\end{theorem}
\begin{remark}
\source{rigid-local-systems-multiplicative-eigenvalue-problem}
\notready\label{added}
\source{rigid-local-systems-multiplicative-eigenvalue-problem}
Under the above correspondence,
F-line bundles on $\Par_{n,\mathcal{O},S}$ (and the corresponding F-vertices) on the (B) side are therefore in strange duality with the unitary rigid local systems on (A) side.
\end{remark}
\subsection{Proof of Theorem \ref{christmas}}
We first prove the  contrapositive of  the reverse direction in (i): Suppose,
\begin{equation}\label{decompo}
\source{rigid-local-systems-multiplicative-eigenvalue-problem}
\mathcal{B}(\vec{\lambda},\ell)=\mathcal{B}(\vec{\lambda'},\ell')\otimes \mathcal{B}(\vec{\lambda''},\ell'')
\end{equation}
with $\vec{\lambda}=\vec{\lambda'}+\vec{\lambda''}$ and $\ell=\ell'+\ell''$, and $\vec{\lambda},\vec{\lambda'}$ and $\vec{\lambda''}$ are vectors of  normalised weights, i.e., their $n$th rows are empty.   Assume the two factors $\mathcal{B}(\vec{\lambda'},\ell')$ and $\mathcal{B}(\vec{\lambda''},\ell'')$ both have non-zero global sections. They then are necessarily of grade zero. The correspondence (A) to (B) for these factors give two (families) of unitary local systems $\mt'$ and $\mt''$ of ranks $\ell'$ and $\ell''$ respectively. We claim that $\mt'\oplus\mt''$ is a local system of rank $\ell$, with local monodromy $\mathcal{A}$: Write $\lambda'^i=\sum_{b=1}^{n-1} c'^b_i\omega_b$ and $\lambda'^i=\sum_{b=1}^{n-1} c''^b_i\omega_b$. Then $\bar{A}_i$ has eigenvalues $\exp(2\pi \sqrt{-1}b/n)$ repeated $c^b_i=c'^b_i+c''^b_i$ times, as desired.

For the other direction of (i), assume that a point in the moduli space corresponding to the strange dual of $\ml=\mathcal{B}(\vec{\lambda},\ell)$ (see Definition \ref{Shosta})  breaks up as a direct sum $\mathcal{T}=\mathcal{T}'\oplus \mathcal{T}''$ of semistable bundles of ranks $\ell'$ and $\ell''$ respectively. Under the correspondence of (A) and (B), suppose $\mt'$ and $\mt''$ correspond to $\mathcal{B}(\vec{\lambda'}, \ell')$  and $\mathcal{B}(\vec{\lambda''}, \ell'')$ respectively. It is now easy to see that
\eqref{decompo} holds on $\Par_{n,\mathcal{O},S}$, and the two factors are effective by Proposition  \ref{schubertt}.


If $\ml=\mathcal{B}(\vec{\lambda},\ell)$ is an F-line bundle on   $\Par_{n,\mathcal{O},S}$, then $\ml'=\mathcal{B}(\vec{\mu},n)$ on the moduli stack  $\Par_{\ell,\widetilde{\mathcal{N}},S}$ has a one dimensional space of sections (see equation \ref{GepnerPP}), making the corresponding moduli space, of representations of the fundamental group with local monodromy data $\ma$, a point, i.e., the corresponding representation is rigid as a unitary representation. Since it is also irreducible by (i), it is an irreducible rigid local system (see Remark \ref{RigidRem}).

If the strange dual gives a rigid unitary representation, then it obviously does not break up as a direct sum of semistable bundles anywhere in its moduli space, and the line bundle has exactly one section (up to scalars). It then follows from (i) that its  strange dual (i.e., $\mathcal{B}(\vec{\lambda},\ell)$) is a Hilbert basis element which has exactly one section up to scalars (use \eqref{GepnerPP}). Writing $\mathcal{B}(\vec{\lambda},\ell)=\mathcal{O}(E)$ we see that $E$ is irreducible and not a multiple since it is a Hilbert basis element. Therefore, $\mathcal{B}(\vec{\lambda},\ell)$ is an F-line bundle, as desired.
\subsubsection{Proof of Theorem \ref{egregium}}
The association of F-line bundles on $\Par_{n,\mathcal{O},S}$ and F-vertices of $P_n(s)$ is a bijection by Proposition \ref{corresp}. Note that this does not use 
strange duality, and follows from the Mehta-Seshadri theorem, geometric invariant theory, and the saturation conjecture (see Section \ref{uber}). 

Theorem \ref{egregium}  therefore reduces to the statement  that $F$-line bundles on $\Par_{n,\mathcal{O},S}$ of level $\ell$ are in one-one bijection with the set of local monodromy data $\ma$ (as in (A)) which correspond to unitary, irreducible, rigid local systems. Now  by Proposition \ref{schubertt}, effective level $\ell$  line bundles on $\Par_{n,\mathcal{O},S}$ are in one-one correspondence with local monodromy data $\ma$ (as in (A)) which correspond to  rank $\ell$ unitary local systems. 


Theorem \ref{egregium} now follows from Theorem \ref{christmas} (ii). Reviewing the proof of Theorem \ref{christmas}, we see that if we are able to 
decompose an effective  line bundle on $\Par_{n,\mathcal{O},S}$ as tensor product of effective line bundles, then on the strange dual side, we can create a reducible representation with monodromy data $\mathcal{A}$ (and vice versa).


\begin{remark}
\source{rigid-local-systems-multiplicative-eigenvalue-problem}
\notready\label{numerics}
\source{rigid-local-systems-multiplicative-eigenvalue-problem}
If $\rigid$  is an irreducible rigid local system of rank $\ell$  on $\Bbb{P}^1-S$ with local monodromies $A_i$ (not necessarily semisimple), then one has \cite[Theorem 1.1.2]{Katz}
\begin{equation}\label{eqRi}
\source{rigid-local-systems-multiplicative-eigenvalue-problem}
\sum_{i=1}^s \dim Z(A_i)= (s-2) \ell^2+2
\end{equation}
Here $Z(A_i)$ is the subgroup of $\op{GL}(\ell,\Bbb{C})$ of matrices that commute with $A_i$. Also note that if $A$ is a semisimple matrix  with eigenvalues $a_1,\dots,a_p$ of multiplicities $n_1,\dots,n_p$, then $\dim Z(A)=\sum_{j=1}^p n_j^2$.
\end{remark}
\begin{lemma}
\source{rigid-local-systems-multiplicative-eigenvalue-problem}
\notready\label{RigidRem}
\source{rigid-local-systems-multiplicative-eigenvalue-problem}
Suppose  $\rigid$ is a  unitary local system on $\Bbb{P}^1-S$. Let $\rho: \pi_1(\Bbb{P}^1-S,b)\to \op{U}(\ell)$ be the  corresponding representation. Then the following conditions on $\rigid$ are equivalent :
\begin{enumerate}
\item[(a)]$\rigid$ is an irreducible rigid local system in the sense of Katz \cite{Katz}.
\item[(b)] $\rigid$ is an irreducible local system, which is rigid as unitary local system (i.e., any $\rho': \pi_1(\Bbb{P}^1-S,b)\to \op{U}(\ell)$ with the same local monodromies as $\rho$ is conjugate to it).
\end{enumerate}
\end{lemma}
\begin{proof}
This follows from the fact that the complexification of the Lie algebra of $\op{U}(\ell)$ is the Lie algebra of $\op{GL}(\ell)$, and the description of tangent spaces of representation varieties. Here is another direct numerical explanation:
To see this, note that the irreducibility requirements in (a) and (b) are the same. Therefore let us assume these to hold. In this situation (a) is equivalent to \eqref{eqRi}  with $A_i=\rho(\gamma_i)$.
Assuming irreducibility, we have a stable point in the Mehta-Seshadri moduli space of $\rigid$ (call this moduli $\mathcal{M}$). Using \cite[Lemma 5.4]{BTIFR}, (b) is equivalent to the same numerical condition:
Let $\rigid\in\mathcal{M}$ correspond to a parabolic bundle whose underlying bundle is $\mw$. We may assume $\mw$ to be generic, and automorphisms of $\mw$ is a group of dimension $\ell^2$ \cite[Lemma 9.1]{BTIFR}, and (b) says that it has an open orbit in the  corresponding variety of partial flags in the fibers of $\mw$ at points of $S$. The corresponding partial  flag variety at the point $p_i\in S$ should be shown to have dimension $\frac{1}{2}(\ell^2-Z(A_i))$. Using Remark
\ref{numerics}, this follows from \cite[Lemma 1.2]{BTIFR}.
\end{proof}

\begin{remark}
\source{rigid-local-systems-multiplicative-eigenvalue-problem}
\notready\label{Hilberty} Let $n$ and $|S|$ be fixed.  There is a constant $\tilde{C}(|S|,n)$ with the following property:
\source{rigid-local-systems-multiplicative-eigenvalue-problem}
 Suppose $\mathcal{T}$ is an irreducible  rank $\ell$ unitary local system (possibly non-rigid) on $\Bbb{P}^1-S$ with local monodromy data $\mathcal{A}=(\bar{A}_1,\dots,\bar{A}_s)$ with $\bar{A}_i^n=I_{\ell}$, such that any unitary local system with local monodromy data $\mathcal{A}$ is irreducible. Then, $\ell\leq \tilde{C}(|S|,n)$. This boundedness property follows from part (i) of Theorem \ref{christmas}, and the finiteness  of the number of elements in the Hilbert basis of a saturated semigroup.

Local systems  $\mathcal{T}$ with the  above property are not necessarily rigid. By Theorem \ref{christmas} the strange duals of non F-vertices of $P_n(s)$ give examples of families of such $\mt$ which are not rigid.  An example of a non F-vertex in $P_9(3)$ is given in Section \ref{Rex}. The corresponding strange dual is an one dimensional family of such unitary $\mt$ of rank $3$, with local monodromies $(\zeta^3,\zeta^7,\zeta^8)$ at $p_1$, and ($1,\zeta^3,\zeta^6)$ at $p_2$ and $p_3$, with $\zeta=\exp(2\pi\sqrt{-1}/9)$.  We are unable to construct any such $\mt$ explicitly (so this is only an existence statement).

We find reducible but non-unitary $\mt$ with the same local monodromy data: of the form $\mt_1\oplus\mt_2$, where $\mt_1$ is rank one with local monodromies $\zeta^3$ at all three points, and $\mt_2$ of rank $2$ with local monodromies $(\zeta^7,\zeta^8)$, $(1,\zeta^6)$ and $(1,\zeta^6)$. Such a $\mt_2$ arises from a hypergeometric system with $\alpha_1=7/9$, $\alpha_2=8/9$, $\beta_1=0$ and $\beta_2=3/9$. $\mt_2$ is not unitary (see Theorem \ref{BHhere}).

\end{remark}
\begin{remark}
\source{rigid-local-systems-multiplicative-eigenvalue-problem}
\notready\label{stretch2}
\source{rigid-local-systems-multiplicative-eigenvalue-problem}
In Theorem \ref{christmas}, the hypotheses on the (moduli of) unitary local systems coming from the strange duals are invariant under scaling the choice of $n$, i.e., if local monodromies are $n$th roots of unity, they are also $(nm)$th roots of unity for any positive integer $m$. Under strange duality, scaling is strange dual to stretching: The $m$-stretch of a line bundle
$\mathcal{B}(\vec{\lambda},\ell)$ on  $\Pic^+(\Par_{n,\mathcal{O},S})$ is a line bundle $\mathcal{B}(\vec{\lambda}[m],\ell)$ on  $\Pic^+(\Par_{mn,\mathcal{O},S})$. Here if $\lambda=\sum_{i=1}^{n-1}  c_i\omega_i$ is a weight for $\mathfrak{sl}_n$ then $\lambda[m] = \sum_{i=1}^{n-1}  c_i\omega_{mi}$ is a weight for $\mathfrak{sl}_{mn}$  (see \cite[Prop 1.3]{GG} for an appearance of this operation in the context of Chern classes of conformal blocks).

We can also scale in $\ell$ on the (B)-side in Proposition \ref{schubertt}, replacing $\mathcal{B}(\vec{\lambda},\ell)$ by a power $\mathcal{B}(m\vec{\lambda},m\ell)$, but now (A) has rank
$m\ell$. The cone structure on the (B)-side in $\Pic^+(\Par_{n,\mathcal{O},S})$ therefore goes into an operation corresponding to direct sums on the (A)-side.

According to Theorem \ref{christmas}, the properties of being a Hilbert basis element, and a F-line bundle (hence for F-vertices), are invariant under stretching. This invariance property for F-vertices (similarly for all vertices), under stretching,  was proved previously in unpublished joint work with  Gibney and Kazanova using the natural (``direct sum'') map $\Par_{n,\mathcal{N},S}\to \Par_{mn,\mathcal{N}^m,S}$, quantum saturation and Fulton conjectures (Proposition   \ref{qTheorems}), and strange duality \ref{BWag}.
\end{remark}

\subsection{Proof of Corollary \ref{egregiumprime}}
\begin{defi}
\source{rigid-local-systems-multiplicative-eigenvalue-problem}
\notready\label{aliB}
\source{rigid-local-systems-multiplicative-eigenvalue-problem}
If $\lambda=\sum_{a=1}^{n-1} c_a\omega_a$ is a dominant integral weight of $\op{SL}(n)$, and $\op{gcd}(m,n)=1$, let $T_m(\lambda)$ be the dominant integral weight $\sum c_a\omega_{ma \pmod{n}}$.
\end{defi}

\begin{lemma}
\source{rigid-local-systems-multiplicative-eigenvalue-problem}
\notready\label{calcul}
\source{rigid-local-systems-multiplicative-eigenvalue-problem}
Suppose $\sigma\in \op{Gal}(\bar{\Bbb{Q}}/\Bbb{Q})$ has image $\bar{m}\in (\Bbb{Z}/n\Bbb{Z})^*=\op{Gal}(\Bbb{Q}(\zeta_n)/\Bbb{Q})$.

\begin{enumerate}
\item Then $v(\sigma(\mathcal{A}))= T_m(v(\mathcal{A}))$ ($T_m$ is as defined in Definition  \ref{aliA}).
\item  $\mathcal{B}(\vec{\lambda},\ell)$ is indivisible as a grade zero line bundle if and only if
$\mathcal{B}(T_m(\vec{\lambda}),\ell))$ is indivisible as a grade zero line bundle ($T_m$ is  defined in Definition  \ref{aliB}).
\end{enumerate}
\end{lemma}
\begin{proof}
Consider a  semisimple class $\bar{A}\in \op{GL}(\ell,\Bbb{C})$ with eigenvalues $\exp(2\pi\sqrt{-1}\mu_j/n)$ with $n>\mu^i_1\geq \mu^i_2\geq \dots \geq \mu^i_{\ell}\geq 0$. Let $\lambda=\mu^T$ be expressed in the form $\sum_{a=1}^{n-1} c_a\omega_a$. The eigenvalues of $\bar{A}$ are $\zeta_n^a$ with multiplicity  $c_a$  for $1\leq a \leq n$, and $1$ with multiplicity $\ell-\sum c_a$. The action of $\sigma$ on $\bar{A}$ changes $\lambda$ to $\sum_{a=1}^{n-1} c_a\omega_{ma \pmod n}$ since the element of $\op{Gal}(\Bbb{Q}(\zeta_n)/\Bbb{Q})$ corresponding to $m$ takes $\zeta_n$ to $\zeta_n^m$.
The lemma now follows from a small computation.
\end{proof}
We now prove Corollary \ref{egregiumprime}. Suppose there is a rigid local system $\rigid$  with local monodromy data $\mathcal{A}$, with finite global monodromy. Suppose the corresponding matrices are $(A_1,\dots,A_s)$. Then we may assume
$A_i$ have coefficients in $\bar{\Bbb{Q}}$ (in fact $\rigid$ is defined over the ring of integers of a cyclotomic
number field by a theorem of Katz \cite{Katz}). Then the twists $(\sigma(A_1),\dots,\sigma(A_s))$ also have finite monodromy for all $\sigma\in \op{Gal}(\bar{\Bbb{Q}}/\Bbb{Q})$. Hence there is a unitary irreducible rigid local system corresponding to $\mathcal{A}^{\sigma}$ for every $\sigma$. We can complete the forward implication in the corollary using Theorem \ref{egregium} and Lemma \ref{calcul}.

For the reverse implication let $\rigid$ be the unitary irreducible rigid local system corresponding to $\mathcal{A}$. By rigidity, and our assumptions,  all Galois conjugates of $\rigid$ are unitary. Hence by a standard argument, see e.g., \cite[Theorem 11]{BLocal}, Katz's theorem that $\rigid$ is
defined over the ring of integers of a cyclotomic number field implies that $\rigid$ has finite global monodromy.
\subsection{More consequences and examples}
Theorem \ref{christmas} therefore imposes numerical conditions on the weights of F-line bundles:
\begin{corollary}
\source{rigid-local-systems-multiplicative-eigenvalue-problem}
\notready\label{katzkatz}
\source{rigid-local-systems-multiplicative-eigenvalue-problem}
Suppose $\mathcal{B}(\vec{\lambda},\ell)$ is an F-line bundle on $\Par_{n,\mathcal{O},S}$. Write $\lambda^i=\sum_{b=1}^{n-1}c^b_i\omega_b$ for $i=1,\dots,s$. Then,
\begin{equation}\label{egale}
\source{rigid-local-systems-multiplicative-eigenvalue-problem}
\sum_{i=1}^s \bigl((\ell-\sum_{b=1}^{n-1} c^b_i)^2 \ + \ \sum_{b=1}^{n-1}(c^b_i)^2\bigr) =(s-2)\ell^2+2.
\end{equation}
\end{corollary}
\begin{proof}
Then the transpose of $\lambda^i$ has rows of length $b$ repeated $c^b_i$ times and a row of length $0$ repeated $\ell-\sum_{b=1}^{n-1} c^b_i$ times. Equation \eqref{egale} now follows from Theorem \ref{christmas} and equation \eqref{eqRi}.
\end{proof}










\iffalse

We now have the following corollary of Theorem \ref{christmas}(ii), and Proposition \ref{corresp},
\begin{corollary}
\source{rigid-local-systems-multiplicative-eigenvalue-problem}
\notready\label{bounded}
\source{rigid-local-systems-multiplicative-eigenvalue-problem}
The following are equivalent
\begin{enumerate}
\item There is a rigid unitary representation with  local monodromy data $\mathcal{A}$.
\item $\mathcal{B}(\vec{\lambda},\ell)$ is a F-line bundle on $\Par_{n,{\mathcal{O}},S}$.
\end{enumerate}
When these conditions are satisfied, $v(\ma)$ is a F-vertex of $P_n(s)$ (also see Lemma \ref{corresp}).

The association of F-line bundles and F-vertices is bijective (by Proposition \ref{corresp}). Since the polytope $P_n(s)$ has only finitely many vertices, the set of  unitary irreducible rigid local systems on $\Bbb{P}^1-S$ whose local monodromies are $n$th
roots of unity, is a finite set, i.e., the ranks $\ell$ of such local system is bounded above by a constant $C(|S|,n)$.
\end{corollary}

The set of ranks of rigid local systems of finite global monodromy on $\Bbb{P}^1-S$ whose local monodromies are $n$th roots of unity, is also therefore bounded above by the same constant $C(|S|, n)$, and hence there are only finitely many irreducible rigid local systems on $\Bbb{P}^1-S$ whose local monodromies are $n$th roots of unity, and have finite global monodromy.
\fi

\subsection{Some examples}\label{smalln}
\source{rigid-local-systems-multiplicative-eigenvalue-problem}
\iffalse
\begin{example}
\source{rigid-local-systems-multiplicative-eigenvalue-problem}
\notready\label{identical}
\source{rigid-local-systems-multiplicative-eigenvalue-problem}
We compute the strange dual (see Definition \ref{Shosta}) of the F-line bundle from Section \ref{oldie}. Recall that we had an F-line bundle $\mathcal{B}(\vec{\lambda},\ell)$ on $\Par_{4,\mathcal{O},S}$ with $S=\{p_1,p_2,p_3\}$ with $\ell=2$ and $\lambda^1=\lambda^2=\lambda^3=\omega_1+\omega_3$. It is easy to see that Equation \eqref{egale} is valid: $2= (2^2)(2-3)+ (3) (1^2+1^2))$.

The strange dual is a line bundle $\mathcal{B}(\vec{\lambda}^T, 4)$ on $\Par_{4,\mathcal{O}(-1),S}$ where the transposes are taken in a $4\times 2$ box. It is easy to see that the corresponding matrix equation (see Corollary \ref{sansD}) in the unitary group $U(2)$ is $A_1A_2A_3=C$ where $A_1,A_2,A_3$ are conjugate to the diagonal matrix with  eigenvalues $\exp(2\pi\sqrt{-1}/4)$ and $\exp(6\pi\sqrt{-1}/4)$, i.e., $\sqrt{-1}$  and $-\sqrt{-1}$, and $C=I_2$. This is a rigid unitary local system with finite global monodromy, since all Galois conjugates are identical (this local system with finite monodromy is the simplest from Schwarz' list \cite{Schwarz}, corresponding to parameters $(1/2,1/2,1/2)$).
\end{example}
\fi
In this section we use known explicit descriptions of vertices $P_n(3)$ for $n\leq 6$ to list all unitary irreducible rigid local systems whose local monodromies are $n$th roots of identity.  To do this we  start with a F-vertex use Proposition \ref{corresp} and Remark \ref{inpractice} to find the corresponding F-line bundle, and then apply Theorem \ref{christmas} to find the corresponding rigid local system. For $n\leq 6$ all vertices of $P_n(3)$ are F-vertices.


By direct computation $P_2(3)\subseteq[0,\frac{1}{2}]^3$ is a polytope with vertices $(\frac{1}{2}, \frac{1}{2},0)$ and various permutations. Note that $\Delta_2=\{(a,-a): a\geq 0,\ a\leq -a +1\}=[0,\frac{1}{2}]$.
The strange duals of these vertices are of level one, so we only get rigid local systems of rank $1$. Therefore there are no unitary and irreducible  rigid local systems on $\Bbb{P}^{1}-\{p_1,p_2,p_3\}$ of rank $>1$ whose local monodromies have eigenvalues $\pm 1$.

By an explicit computation of $P_3(3)$ a similar conclusion holds for $n=3$: There are no unitary and irreducible  rigid local systems on $\Bbb{P}^{1}-\{p_1,p_2,p_3\}$ of rank $>1$ whose local monodromies have eigenvalues third roots of unity.

For $n=4$, by \cite[Section 7]{BLocal}, there are F-vertices of levels are bigger than one: These correspond to the  F-line bundle $\mathcal{B}(\vec{\lambda},\ell)$ with $\ell=2$ and $\lambda^1=\lambda^2=\lambda^3=\omega_1+\omega_3$. Therefore in the corresponding the rank $2$ local system   all local monodromies are conjugate to the diagonal matrix with entries $\exp(2\pi\sqrt{-1}/4)$ and $\exp(6\pi\sqrt{-1}/4)$, i.e., $\sqrt{-1}$ and $-\sqrt{-1}$.  The corresponding local system corresponds to the equation $\sigma_1\sigma_2\sigma_3=I_2$ where
\begin{equation}\label{dfnsigma}
\source{rigid-local-systems-multiplicative-eigenvalue-problem}
   \sigma_1=
  \left[ {\begin{array}{cc}
   0 & 1 \\
   -1 & 0 \\
  \end{array} }\right],   \sigma_2=
  \left[ {\begin{array}{cc}
   -\sqrt{-1} & 0 \\
   0 & \sqrt{-1} \\
  \end{array} }\right],  \sigma_3=
  \left[ {\begin{array}{cc}
   0 & \sqrt{-1} \\
   \sqrt{-1} & 0 \\
  \end{array} }\right].
  \end{equation}
This local system and all of its twists by $4$th roots of unity  are the only unitary irreducible rigid local systems on $\Bbb{P}^{1}-\{p_1,p_2,p_3\}$ of rank $>1$ whose local monodromies have eigenvalues fourth roots of unity.
(By twisting we mean the following. Let $a_1,a_2,a_3$ be $n$th roots of unity with $a_1a_2a_3=1$.  Then the twist is just the tensor product with the rank one local system given by $(a_1,a_2,a_3)$). They all have finite global monodromy by Corollary \ref{egregiumprime}, and appear in Schwarz's list from the 19th century \cite{Schwarz}, and are of hypergeometric type (up to a twist).

For $n=5,6$ we use computations of $P_n(3)$ reported in \cite{thaddy}.

For $n=5$ there are two orbits under $\tau_5(3)\rtimes S_3$ (see Section \ref{ABWin})  of vertices of $P_5(3)$, both consisting of F-vertices. The first one is the orbit of $(0,0,0,0,0)^3$ whose strange duals are rank one. The second is the orbit of $(\frac{1}{2},0,0,0,-\frac{1}{2})^3$, the point $(\frac{1}{2},0,0,0,-\frac{1}{2})$ corresponds to $(1/2,0,0,1/2)\in\Delta'_5$ in the notation of Remark \ref{aliA}.
The corresponding unitary irreducible rigid local system is of rank $2$. By Theorem \ref{KatzQ}, none of the second kind of local systems have finite global monodromy (which can also be checked here by using Corollary \ref{egregiumprime}).

\iffalse
There are four choices of $m$: $m=1, m=4$ give the original point and its (ordinary) dual. Therefore the only check for finiteness of global monodromy is for $m=2$, i.e., whether $(0,1/2,1/2,0)^3\in(\Delta'_5)^3$ (which corresponds to $(1/2,1/2,0,-1/2,-1/2)^3\in \Delta_5^3$) is a vertex of $P_5(3)$. It is easy to check that this is not in the $\tau_5(3)\rtimes S_3$  orbit of $(\frac{1}{2},0,0,0,-\frac{1}{2})^3$: Therefore there are no irreducible  rigid local systems on $\Bbb{P}^{1}-\{p_1,p_2,p_3\}$ of rank $>1$, with finite global monodromy, whose local monodromies have eigenvalues fifth roots of unity.
\fi


For $n=6$, let us first note that the Galois group $\op{Gal}(\Bbb{Q}(\zeta_n)/\Bbb{Q})$ has order two, and hence we just need to list the strange duals of all vertices of $P_6(3)$, all of these will have finite global monodromy (the action on $\Delta_n^3$ is either the identity or by ordinary duals, and both preserve $P_n(3)$). There are three orbits under $\tau_6(3)\rtimes S_3$ (see Section \ref{ABWin})  of vertices of $P_5(3)$,  all consisting of F-vertices. There strange duals give us the following local systems
on $\Bbb{P}^{1}-\{p_1,p_2,p_3\}$. These and their central twists by $6$th roots of unity are the only unitary (hence finite for $n=6$) rigid local systems on $\Bbb{P}^{1}-\{p_1,p_2,p_3\}$ with local monodromies  sixth roots of unity, here $\zeta=\exp(\frac{2\pi\sqrt{-1}}{6})$. They are all of classical
hypergeometric type (up to a twist)  as in Section \ref{classH}. We indicate the classical type below.
\begin{enumerate}
\item The first local system is a rank $2$ local system with  all local monodromies conjugate, the eigenvalues are $(\zeta^5,\zeta)$. The twist consisting of local monodromies $(\zeta^5,\zeta)$, $(1,\zeta^2)$, and $ (1,\zeta^4)$ corresponds to a rank $2$ hypergeometric system with $\alpha_1=\frac{1}{6}$, $\alpha_2=\frac{5}{6}$, $\beta_1=0$ and  $\beta_2=4/6$.
\item The second local system is a rank $2$ local system with  eigenvalues of local monodromies  given by
$(\zeta,\zeta^5), (1,\zeta^3)$, and $(1,\zeta^3)$.  This corresponds to a rank $2$ hypergeometric system with $\alpha_1=\frac{1}{6}$, $\alpha_2=\frac{5}{6}$, $\beta_1=0$ and  $\beta_2=3/6$.
\item The third local system is a rank $3$ local system with eigenvalues of local monodromies $(1,1,\zeta^3)$, $(1,\zeta^2,\zeta^4)$, and $(\zeta,\zeta^3,\zeta^5)$. This  corresponds to a rank $3$ hypergeometric system with $\alpha_1=\frac{1}{6}$, $\alpha_2=\frac{3}{6}$,  and $\alpha_3= \frac{5}{6}$, $\beta_1=0$,  $\beta_2=\frac{2}{6}$ and $\beta_3=\frac{4}{6}$.
\end{enumerate}

\subsection{Special properties of rigid unitary local systems and the proof of Theorem \ref{KatzQ}} \label{KatzSection}
\source{rigid-local-systems-multiplicative-eigenvalue-problem}
\begin{proposition}
\source{rigid-local-systems-multiplicative-eigenvalue-problem}
\notready\label{propP}
\source{rigid-local-systems-multiplicative-eigenvalue-problem}
Let $\mt$ be a rank $\ell$ irreducible rigid local system with semisimple local monodromy data
$\mathcal{A}=(\bar{A}_1,\dots,\bar{A}_s)$ in $\op{GL}(\ell,\Bbb{C})$, with $\bar{A}_i^n=1$. Let $i\in[s]$ and let $\alpha$ and $\beta$ be two eigenvalues of $\bar{A}_i$.
\begin{enumerate}
\item If $\mt$ is unitary, then $\alpha^{-1}\beta\neq \zeta_n$ where $\zeta_n=e^{2\pi\sqrt{-1}/n}$ (i.e., two different eigenvalues of the local monodromy at a point of $S$ cannot be ``as close as possible" for unitary rigid irreducible local systems).
\item If $\mt$ has finite global monodromy, then $\alpha^{-1}\beta$ is not a primitive $n$th root of unity.
\end{enumerate}
\end{proposition}
Theorem \ref{KatzQ}, and property (P) from the introduction are clearly  immediate consequences of the above proposition.
\begin{proof}
The Galois group $\op{Gal}(\Bbb{Q}(\zeta_n)/\Bbb{Q})$ acts transitively on the set of primitive $n$th roots of unity and therefore the second part follows from the first part (use that any Galois conjugate of $\mt$ is a unitary local system in addition to being rigid and irreducible).

For the first part we proceed as follows. Assume the contrary, then by the action of $\tau_n(s)$, we may assume $\alpha=1$, $\beta= \exp(2\pi \sqrt{-1}/n)$. Let the corresponding F-line bundle (by the correspondence between data  (A) and (B) in Section \ref{SD}) on $\Par_{n,\mathcal{O},S}$ be $\ml=\mb(\vec{\lambda},\ell)$. Write  $\lambda^i=\sum_{b=1}^{n-1}c^b_i\omega_b$ for $i=1,\dots,s$. We also have from Lemma \ref{divisorE}, an expression $\ml=\mathcal{O}(E)$ where  $E=C(d,r,\mathcal{O},n,\vec{J})$.

By Proposition \ref{enumerative}, parts (C) and (D), and the strange duality operation (see Definition \ref{Shosta}) the following properties hold which imply that $1$ and $\zeta_n$ cannot both be eigenvalues of $A_i$:
\begin{enumerate}
\item If  $1\leq b<n$,  then $\exp(2\pi\sqrt{-1}b/n)$ is an eigenvalue of $A_i$ with multiplicity $c^b_i$ (possibly zero). If this multiplicity is non-zero, then $b\in J^i$ and $b+1\not\in J^i$.
\item The multiplicity of the eigenvalue  $1$ of $A_i$ is  $\ell-\sum_{b=1}^{n-1}c^b_i$ (possibly zero). If this multiplicity is non-zero,  then $n\in J^i$ and $1\not\in J^i$.
\end{enumerate}
\end{proof}
\begin{remark}
\source{rigid-local-systems-multiplicative-eigenvalue-problem}
\notready\label{moreR} The assumption of rigidity is needed in property (P) and Proposition \ref{propP}: Enlarge $S$ by two points and introduce local monodromies at the points that are inverses of each other.
\source{rigid-local-systems-multiplicative-eigenvalue-problem}
\end{remark}


\subsection{Rigid local systems from quantum Schubert calculus}\label{qSC}
\source{rigid-local-systems-multiplicative-eigenvalue-problem}
Theorems \ref{Adagio} and \ref{brahms} produce F-line bundles $\mathcal{O}(D(a,j))$, and hence rigid irreducible unitary local systems as in Corollary \ref{existence}. Let us recall the steps:
\begin{enumerate}
\item [(i)]Start with a situation where
 $\langle\sigma_{I^1}, \sigma_{I^2},\dots,\sigma_{I^s}\rangle_d= 1$ in a Grassmannian $\Gr(r,n)$.
\item[(ii)] Choose $(a,j)$: Let $(a,j)$ be a pair with $a\in[n]$ and $1\leq j\leq s$ such that
\begin{enumerate}
\item $a>1$, $a\in I^j$, and $a-1\not\in I^j$, or
\item $a=1$, $1\in I^j$, $n\not\in I^j$.
\end{enumerate}
\item[(iii)] Define subsets $J^1,\dots,J^s$ of $[n]$ each of cardinality $r$ and an integer $d'$ as follows:
In case (1) let $J^k=I^k$ if $k\neq j$, $J^j=(I^j-\{a\}) \cup \{a-1\}$ and $d'=d$, and in case (2) let  $J^k=I^k$ if $k\neq j$, $J^j=(I^j-\{1\}) \cup \{n\}$ and $d'=d-1$.
\item[(iv)] Construct the F-line bundle $\mathcal{B}(\vec{\lambda},\ell)$ on $\Par_{n,\mathcal{O},n}$
by the formulas $\lambda^i=\sum_{b=1}^{n-1}c^i_b\omega_b$ with formulas for $\ell$ and $c^i_b$
as given in Theorem \ref{brahms}.
 \item[(v)] Finally the rigid local system $\mt$ is of rank $\ell$ on $\Bbb{P}^1-S$ with local monodromy at $p_i$ for $i\in[s]$
 conjugate to a diagonal matrix with eigenvalues $\exp(2\pi\sqrt{-1}b/n)$ with multiplicity $c^i_b$, for $b=1,\dots n-1$, and eigenvalue $1$ with multiplicity $\ell-\sum_{b=0}^{n-1} c^i_b$. Note that the KZ differential equation which gives rise (up to a cyclic twist) to $\mathcal{T}$ is given explicitly by Theorem \ref{KZe} and Proposition \ref{generalKZ}, since the data of the divisor $E$ in the proof of Theorem \ref{KZe} is explicit in this situation.
 \end{enumerate}

 Quantum Schubert calculus gives a method of computing all $\langle\sigma_{I^1}, \sigma_{I^2},\dots,\sigma_{I^s}\rangle_d$, and in the case $d=0$ there are many other methods for computing these numbers (Littlewood-Richardson rule, the honeycomb and puzzle formulations of Knutson and Tao). As for systematic methods of obtaining situations where $\langle\sigma_{I^1}, \sigma_{I^2},\dots,\sigma_{I^s}\rangle_d= 1$, we note the following:
\begin{enumerate}
\item Let $\lambda,\mu$ and $\nu$ be dominant integral weights of $\op{SL}_r$ such that there is a Weyl group  element $w\in S_r$ with $\lambda+w\mu=\nu$. In this case, by a result of Kostant \cite[Lemma 4.1]{Kostant}, $V_{\nu}$ appears with multiplicity one in $V_{\lambda}\tensor V_{\mu}$. Since Littlewood-Richardson numbers also compute structure coefficients in the classical cohomology of suitable Grassmannians $\Gr(r,n)$, this result of Kostant produces a family of $\langle\sigma_{I^1}, \sigma_{I^2},\dots,\sigma_{I^s}\rangle_d=1$ with $d=0$.
\item The coefficients that appear in the quantum Pieri rule are all one \cite{bertie}. These correspond to $s=3$ and $\lambda(I^1)=(a,0,\dots,0)$ (or their duals). Some of the unitary local systems that arise from this situation seem to be a subset of  those coming from the classical  hypergeometric systems (Section \ref{hyper}), but there is an infinite series of new ones  as shown by Kiers and Orelowitz (see Section \ref{KO} for their examples).
\item Let $V_1,\dots,V_s$ be  subspaces in general position of a vector space $W$ so that $\sum_{i=1}^s \dim V_i= \dim W+1$. Now there is a unique subspace $T\subset W$ of dimension $s-1$ that meets all $V_i$ non-trivially.
This sets up a classical intersection number $1$ situation in $\Gr(s-1,W)$ (we learned of this example from Nori): $\langle\sigma_{I^1},\dots,\sigma_{I^s}\rangle_0=1$ where $\lambda(I^j)=(a_j,0,\dots,0)$ (see Definition  \ref{correspondence}) with $a_j = n-(s-1) +1-\dim V_j$. The corresponding unitary irreducible rigid local systems seem to include some classical Pochhammer systems (Section \ref{pokhie}).
\item There are examples of Hobson \cite{Hobson} of rank one conformal blocks in type A, which then give rise
to a situation where $\langle\sigma_{I^1}, \sigma_{I^2},\dots,\sigma_{I^s}\rangle_d= 1$ using the Witten dictionary (see Section \ref{QCB}).
\item Suppose we have a classical intersection number one   situation in a Grassmannian $\Gr(r,n)$, i.e., $\langle\sigma_{I^1}, \sigma_{I^2},\dots,\sigma_{I^s}\rangle_d= 1$ with $d=0$. We can then create an intersection number one situation in $\Gr(r,n+1)$:   $\langle\sigma_{J^1}, \sigma_{J^2},\dots,\sigma_{J^s},\sigma_{J^{s+1}}\rangle_0= 1$ with $J^{s+1}=\{n-1,n-2,\dots,n-r-1\}$ and $J^i$ obtained by adding $1$ to elements of $I^i$ for $i\in[s]$.  Running the above process from the two intersection one situations above, one gets two local systems (of possibly different ranks), one on $\Bbb{P}^1-\{p_1,\dots,p_s\}$ and the other on $\Bbb{P}^1-\{p_1,\dots,p_s,p_{s+1}\}$, and the local monodromies have eigenvalues that are $n$th and $(n+1)$th roots of unity respectively.
\item Any classical GW intersection number (i.e., $d=D=0$) in a Grassmannian $\Gr(r+1,r+\ell+2)$ breaks up as a sum of two GW numbers for $d=1$ on Grassmannians $\Gr(r,r+\ell+1)$ and $\Gr(\ell+1,r+\ell+1)$ by \cite[Section 9.9]{BGM}. If the classical intersection number is one, then exactly of the two GW numbers is one.
\item Theorem \ref{Adagio} inductively produces situations where $\langle\sigma_{I^1}, \sigma_{I^2},\dots,\sigma_{I^s}\rangle_d= 1$, since the line bundles $D(a,j)$ are F-line bundles: The new Grassmannian
is $\Gr(n,n+\ell)$ where $\ell$ is the level of $\mathcal{O}(D(a,j))$ (and using Lemma \ref{newlabel2}).
\end{enumerate}
In Sections \ref{E1} and \ref{KO}, we list examples of unitary rigid local systems that are not of hypergeometric and Pochhammer type (which have at least one  matrix with only two distinct eigenvalues). Several examples of $\langle\sigma_{I^1}, \sigma_{I^2},\dots,\sigma_{I^s}\rangle_d= 1$ for $n\leq 16$ are given in \cite{OO} which also includes a discussion of computational aspects of Gromov-Witten invariants.
\subsection{All unitary irreducible rigid local systems}\label{qSC2}
\source{rigid-local-systems-multiplicative-eigenvalue-problem}
  One may ask if all unitary irreducible rigid local systems with local monodromy $n$ roots of identity arise from Corollary \ref{existence}, i.e., the procedure described in Section \ref{qSC}. By Theorem \ref{christmas} (ii), this is equivalent to the question  of whether all F-line bundles on $\Par_{n,\mathcal{O},S}$ are of the form $\mathcal{O}(D(a,j))$. By Lemma \ref{powercord} (also see Proposition  \ref{practical}), this is equivalent to asking
\begin{question}
\source{rigid-local-systems-multiplicative-eigenvalue-problem}
\notready\label{Q}
\source{rigid-local-systems-multiplicative-eigenvalue-problem}
Let $\mathcal{B}(\vec{\lambda},\ell)$ be an F-line bundle on $\Par_{n,\mathcal{O},S}$. Write $\lambda^i=\sum_{b=1}^{n -1} c^b_i\omega_b$. Then is it necessarily true that  there is an $i$ such that
$c^b_i=1$ for some $b$, or $\ell -\sum_{b=1}^{n-1} c^b_i=1$?
\end{question}
The strange dual way of asking this question is the following (see the proof of Corollary \ref{katzkatz}): Let $\mathcal{T}$ be a irreducible rigid unitary local system on $\Bbb{P}^1-S$. Is it necessarily the case that the local monodromy of $\mathcal{T}$ around some $p_i\in S$ has an eigenvalue of multiplicity one? See Section \ref{wilson} for an example, due to A. Wilson, in which
$c^b_i\neq 1$ for all $b$ and $i$ (coming from a classical context), but in this example  $\ell -\sum_{b=1}^{n-1} c^b_i=1$ for some $i$. The rigidity equation \eqref{eqRi} does not settle this:
For example a rigid local system of rank $7$ with $|S|=3$, and multiplicity decomposition $(3,2,2)$
at each point is allowed by \eqref{eqRi}. We do not know if there are any unitary irreducible rigid local systems with these eigenvalue decompositions of local monodromies.

Note the hypergeometric and Pochhammer unitary local systems satisfy Question
\ref{Q}, and hence arise from Corollary \ref{existence}.

\subsection{Rigid irreducible local systems without the unitarity condition}\label{afterN}
\source{rigid-local-systems-multiplicative-eigenvalue-problem}
We consider the following question:
\begin{question}
\source{rigid-local-systems-multiplicative-eigenvalue-problem}
\notready Fix a natural number $n$. Consider the set of  irreducible rigid local systems (possibly non-unitary) on $\Bbb{P}^1-S$  such that the local  monodromies are semisimple with eigenvalues $n$ roots of unity.
Is this set finite, i.e., are the ranks $\ell$ of such local systems bounded above?
\end{question}
In the terminology of Katz's book \cite[Section 6.1]{Katz},  the above question asks, in characteristic zero,  whether  realizable  classes in $\op{NumData}(D,\Gamma)$, such that the ``local monodromies'' are semisimple, form a finite set when $\Gamma$ is the cyclic group of order $n$. We do not know the answer to this question.  We examine  below whether plausible (\cite[Definition  6.3.3]{Katz}) elements  in $\op{NumData}(D,\Gamma)$ with semisimple local monodromies form a finite set.


We consider
two natural  necessary conditions on the local monodromy data $\mathcal{A}=(\bar{A}_1,\dots,\bar{A}_s)$ of such a local system of rank $\ell$. The first one is the rigidity equation  \eqref{eqRi}. The second one, a consequence of irreducibility, is the following condition  which can be applied to all cyclic twists of the data (see \cite[Lemma III.9.10]{MM}):
\begin{equation}\label{irreducible1}
\source{rigid-local-systems-multiplicative-eigenvalue-problem}
\sum_{i=1}^s \rk(A_i-I_{\ell})\geq 2\ell
\end{equation}
(This is the inequality that $\chi(\Bbb{P}^1-S, j_*\mathcal{T})\leq 0$ for irreducible $\mathcal{T}$ and $j:\Bbb{P}^1-S\to \Bbb{P}^1$, and follows from Poincare duality, \cite[Section 1.1]{Katz}.)

These two conditions do not seem to put an upper bound on $\ell$ in terms of $n$ by the following argument. Fix $s=3$ for concreteness. We can then put \eqref{irreducible1} in the following form
\begin{equation}\label{irreducible2}
\source{rigid-local-systems-multiplicative-eigenvalue-problem}
\sum_{i=1}^3 n_{i,0}\leq \ell,
\end{equation}
where $n_{i,0}$ is the rank of the zero eigenspace of $A_i$.
Consider ``two extreme cases":
\begin{enumerate}
\item All local monodromies are the same with the multiplicity of each $n$th root of unity the same.
In this case \eqref{irreducible2} is valid  for all cyclic twists since $3(\ell/n)\leq \ell$ if $n\geq 3$. The left hand side of  the equality \eqref{eqRi} would then be $3\ell^2/n$ which is less than $\ell^2+2$ (since $n>3$).
\item There are only two eigenvalues at each of the  three points, with multiplicities $\ell/2$ and $\ell/2$ each. The left hand side of the equality \eqref{eqRi} would then be $3\ell^2/2$ which is greater than $\ell^2+2$ for  $\ell$ large. For \eqref{irreducible2} to fail, all three matrices must have a common eigenvalue (after central twisting). But that does not seem to be forced by the numerical requirements.
\end{enumerate}
Therefore assuming \eqref{irreducible2}, and  that $\ell$ is large  respect to $n$, it is not the case that equality \eqref{eqRi} always fails in the same direction. One could also check whether Katz' algorithm
applied to the data formally produces a rank which is $\geq 0$. When $s=3$, the rank is of the form $2\ell -(\ell_1 +\ell_2+\ell_3)$ where $\ell_i$ are ranks of suitable eigenspaces at the three points (see e.g., the proof of Theorem III.9.5 in \cite{MM}). In both of the extreme cases this number is positive.
